\input{../tex-inputs/latex-leseansicht-vorspann}

\section[Arthur Schnitzler an Gustav Schwarzkopf, 28. 1. 1914]{L04504 Arthur Schnitzler an Gustav Schwarzkopf, 28. 1. 1914}
\nopagebreak\mylabel{L04504v}
\rehead{ }\normalsize\beginnumbering\briefempfaengerindex{Schwarzkopf, Gustav@\textsc{Schwarzkopf, Gustav}!zzzSchnitzler, Arthur@\emph{von Arthur Schnitzler}!1914-01-281@{28. 1. 1914}|(be}
\toendnotes[C]{\smallbreak\pagebreak[2]}
\correspDesc{Versand  durch Arthur Schnitzler am 28. 1. 1914 in München
\newline{}Übermittlung  am 29. 1. 1914 in München
\newline{}Erhalt  durch Gustav Schwarzkopf im Zeitraum [29. 1. 1914
                  – 2. 2. 1914?] in Wien}\toendnotes[C]{\smallbreak}
\Standort{DLA, A:Schnitzler, HS.1985.1.1897.}
\physDesc{Bildpostkarte, 349 Zeichen
\newline{}Handschrift: Bleistift, deutsche Kurrent
\newline{}Versand: Stempel: »\nobreak{}\oindex{München@\textbf{München}|pwk}München, 29. 1. 14., 1–2 N\nobreak{}«.  }\toendnotes[C]{\smallbreak}
\pstart
           \noindent{}\centering{}{\pb}\textcolor{gray}{\textbf{\textbf{München} Zweibrückenstraße\oindex{Zweibrückenstraße@\textbf{Zweibrückenstraße}|pw} mit Isartor\oindex{Isartorplatz@\textbf{Isartorplatz}|pw}}}\pend
           \pstart{}{\pb}\textsc{Hrn. Gustav Schwarzkopf}\pend{}\pstart{}Wien I\oindex{I., Innere Stadt@\textbf{I., Innere Stadt}|pw}\pend{}\pstart{}\textsc{Tiefer Graben 17}\oindex{Wien@\textbf{Wien}!I., Innere Stadt@\textbf{I., Innere Stadt}!Tiefer Graben 17@\textbf{Tiefer Graben 17}|pw}\pend{}{\bigskip}\vspace{1em}
\pstart
           \raggedleft{}28. 1. 914\pend
           \vspace{0.5em}
\pstart
           lieber Guſtav, vielen Dank für die \label{K_L04504-1v}\edtext{Karte}{\lemma{\textnormal{\emph{Karte}}}\Cendnote{\textnormal{XXXX Auszeichnungsfehler: Dokument L04254 nicht gefunden.}}}\label{K_L04504-1}. Am \label{K_L04504-2v}\edtext{Samſtag}{\lemma{\textnormal{\emph{Samstag}}}\Cendnote{\textnormal{Es geht um die Premiere von \emph{Der
                        Kammersänger}\pwindex{Wedekind, Frank 24.\,7.\,1864 Hannover – 9.\,3.\,1918 München@\textsc{Wedekind, Frank} (24.\,7.\,1864 Hannover – 9.\,3.\,1918 München)!Kammersänger. Drei Scenen@\strich\emph{Der Kammersänger. Drei Scenen}|pwk} (Frank Wedekind\pwindex{Wedekind, Frank 24.\,7.\,1864 Hannover – 9.\,3.\,1918 München@\textsc{Wedekind, Frank} (24.\,7.\,1864 Hannover – 9.\,3.\,1918 München)|pwk}),
                        \emph{Boubouroche. Lebensbild in 2 Acten}\pwindex{Courteline, Georges 25.\,6.\,1858 Tours – 25.\,6.\,1929 Paris@\textsc{Courteline, Georges} (25.\,6.\,1858 Tours – 25.\,6.\,1929 Paris)!Boubouroche. Lebensbild in 2 Acten@\strich\emph{Boubouroche. Lebensbild in 2 Acten}|pwk}
                        (Georges Courteline\pwindex{Courteline, Georges 25.\,6.\,1858 Tours – 25.\,6.\,1929 Paris@\textsc{Courteline, Georges} (25.\,6.\,1858 Tours – 25.\,6.\,1929 Paris)|pwk}) und \emph{Literatur}\pwindex{Schnitzler, Arthur 15. 5. 1862 Wien – 21. 10. 1931 ebd.@\textsc{Schnitzler, Arthur} (15. 5. 1862 Wien – 21. 10. 1931 ebd.)!Literatur@\strich\emph{Literatur}|pwk} (Schnitzler)\eventindex{Burgtheater@\textbf{Burgtheater}!Premiere von Der Kammersänger, Boubouroche, Literatur, 30.1.1914@Premiere von Der Kammersänger, Boubouroche, Literatur, 30.1.1914|pwk} am \emph{Burgtheater}\orgindex{Burgtheater@Burgtheater|pwk}.
                     Schwarzkopf\pwindex{Schwarzkopf, Gustav 7.\,11.\,1853 Wien – 13.\,11.\,1939 ebd.@\textsc{Schwarzkopf, Gustav} (7.\,11.\,1853 Wien – 13.\,11.\,1939 ebd.)|pwk} hatte die Hoffnung geäußert,
                     Schnitzler dort zu sehen. Tatsächlich
                  fand die Premiere\eventindex{Burgtheater@\textbf{Burgtheater}!Premiere von Der Kammersänger, Boubouroche, Literatur, 30.1.1914@Premiere von Der Kammersänger, Boubouroche, Literatur, 30.1.1914|pwkv} bereits
                  am Freitag, den 30. 1. 1914 statt. Schnitzler besuchte erst die Vorstellung\eventindex{Burgtheater@\textbf{Burgtheater}!Aufführung von Der Kammersänger, Boubouroche, Literatur, 6.2.1914@Aufführung von Der Kammersänger, Boubouroche, Literatur, 6.2.1914|pwkv} am 6. 2. 1914.}}}\label{K_L04504-2}{ }ſind
               wir keineswegs noch in Wien\oindex{Wien@\textbf{Wien}|pw}, da durch Ofenſetzung
               u dergl meine Zimmer unbewohnbar ſind. \label{K_L04504-3v}\edtext{So{\geminationn}tag oder Montag}{\lemma{\textnormal{\emph{Sonntag oder Montag}}}\Cendnote{\textnormal{Arthur und Olga Schnitzler kehrten am Dienstag, den 3. 2. 1914 nach Wien\oindex{Wien@\textbf{Wien}|pwk} zurück.}}}\label{K_L04504-3} denk ich. Vielleicht halten wir uns ein paar
               Tage \label{K_L04504-4v}\edtext{in Salzburg\oindex{Salzburg@\textbf{Salzburg}|pw}}{\lemma{\textnormal{\emph{in Salzburg}}}\Cendnote{\textnormal{A. S.: \emph{Wiener Schnitzler}, 1. 1. 1914, 2. 1. 1914 und 3. 1. 1914.}}}\label{K_L04504-4} auf. Das \label{K_L04504-5v}\edtext{Concert\eventindex{München@\textbf{München}!Konzert Olga Schnitzler und Vera Schapira, 26.1.1914@Konzert Olga Schnitzler und Vera Schapira, 26.1.1914|pwv}}{\lemma{\textnormal{\emph{Concert}}}\Cendnote{\textnormal{Konzert Olga Schnitzler und Vera Schapira, 26.1.1914. In: A. S.: \emph{Kulturveranstaltungen}, \url{https://schnitzler-kultur.acdh.oeaw.ac.at/pmb183650.html}.}}}\label{K_L04504-5} hier iſt recht {\pb}gut ausgefallen.\pend
           
\pstart
           Wir grüßen Sie herzlichſt.{\\[\baselineskip]}\spacefill\mbox{Arthur}\pend
           \leftskip=0em{}\selectlanguage{ngerman}\endnumbering\briefempfaengerindex{Schwarzkopf, Gustav@\textsc{Schwarzkopf, Gustav}!zzzSchnitzler, Arthur@\emph{von Arthur Schnitzler}!1914-01-281@{28. 1. 1914}|)be}\mylabel{L04504h}
\begin{anhang}
\end{anhang}\newcommand{\dateiname}{L04504}\newcommand{\titel}{Arthur Schnitzler an Gustav Schwarzkopf, 28. 1. 1914}\newcommand{\editorInnen}{Herausgegeben von Jahnke, SelmaMüller, Martin Anton}\input{../tex-inputs/latex-leseansicht-abspann}
